\documentclass[11pt]{article}
\usepackage[T1]{fontenc}
\usepackage[utf8]{inputenc}
\usepackage[a4paper,margin=26mm]{geometry}
\usepackage{amsmath,amssymb,amsthm,booktabs,microtype}
\microtypesetup{expansion=false}
\usepackage[hidelinks]{hyperref}
\newtheorem{theorem}{Theorem}
\newtheorem{proposition}[theorem]{Proposition}
\theoremstyle{remark}\newtheorem{remark}[theorem]{Remark}
\newcommand{\HH}{\mathrm H}
\title{Frozen Light, Two Grooves\\[4pt]\large A note on what Bohm's image does not determine}
\author{Jan Mikul\'ik}
\date{10 November 2026}
\begin{document}
\maketitle

\begin{center}
\emph{A peer view in the spirit of David Bohm.\\
Two friends may ask different questions.
If they witness the same toss, one bit must do for both.}
\end{center}

\medskip
Bohm explained the implicate order with a drop of dye in glycerine.
Stretched, the drop disappears; the motion reversed, it reappears.
The order was put into the whole in advance. He warned that the experiment,
like a hologram, catches only part of the idea: the order is inserted beforehand,
whereas his deeper question was how stable parts arise in a whole.

The finite causal Shannon theorem is not that experiment, and it is not
``only'' a reinterpretation of it. The image says how a prepared whole
unfolds one appearance. The theorem asks how much randomness the whole must
carry so that one decoder serves every admissible appearance at once.
Unfolding alone does not answer that question.

\medskip
\noindent
\textbf{Dictionary, not an identification.}
The record is a finite seed $S$. The frozen rule is a decoder $D(S,g)$.
The needle is the groove $g$, chosen without reading $S$.
The unfolded history is the answer $D(S,g)$.
Nothing in the dictionary is an electron, a tachyon, or a worldline.
Free will is not defined. If the will is the needle, it is chosen after the
record is frozen and it is not a bit of $\HH(S)$. If the will is the draw of
$S$, it is one draw, made before either friend asks. Those are different arrows.

\begin{theorem}[Two grooves]
\label{thm:grooves}
Let each groove be required to return a fair bit.
\begin{enumerate}
\item The shared record $S\in\{0,1\}$ with $D(S,\mathrm{left})=D(S,\mathrm{right})=S$
and $\Pr(S=0)=\Pr(S=1)=1/2$ has $\HH(S)=1$.
Both grooves are fair, and they agree on every atom.
\item The product record $(L,R)$ of two independent fair bits, with
$D((L,R),\mathrm{left})=L$ and $D((L,R),\mathrm{right})=R$, has
$\HH(L,R)=2$. Each groove is fair. The grooves agree with probability $1/2$.
\item The two records induce the same law on each groove separately,
and $\HH(S)=\HH(L,R)/2$. The joint laws differ:
the shared record puts mass $0$ on $(\mathrm{left},\mathrm{right})=(0,1)$.
\item If a finite seed answers the left groove by a deterministic function
of the seed and that answer is fair, then the seed has entropy at least $1$.
The shared record attains the bound.
\end{enumerate}
\end{theorem}

\begin{proof}
A uniform bit has entropy $-\log_2(1/2)=1$.
Four atoms of mass $1/4$ have entropy $2$.
Pushing the shared bit to either groove, or either coordinate of the product,
yields the uniform law on two points.
Data processing gives the lower bound: the entropy of a seed is at least the
entropy of a deterministic function of it, and a fair bit has entropy $1$.
The off-diagonal check is immediate from the supports.
A companion formalization is \texttt{CausalSeed/FrozenLight.lean}.
\end{proof}

The same one bit is the two-friend tree of the note \emph{One Seed, Many Friends}:
one shared fair question, then a private action with deterministic output zero.
There $L(T)=C(T)=1$. An anti-correlated pairing matches each friend's transcript
law and is not the run of any shared decoder. Peer view is that obstruction.
It is not a claim that the friends saw a physical toss.

\begin{remark}
For a finite tree $T$ the companion theorem supplies a seed $S_G$ with
\[
L(T)\le C(T)\le \HH(S_G)\le L(T)+\frac{1+\log_2 e}{2}.
\]
The additive constant is about $1.22135$ bits. Theorem~\ref{thm:grooves}
shows why the image of unfolding cannot be asked to produce this estimate:
two wholes can unfold the same grooves and still differ by a bit.
\end{remark}

What would remain, past the mathematics, is a question nature can refuse.
Matching a replayed history does not identify the record.
Two listeners who can play different grooves can tell the shared toss from
the product. That distinction lives inside the model.
It is not, by itself, a trace in a detector.

\end{document}
